(1) 设左入射压力幅度1、反射 \(r\)，右透射 \(t\)。速度连续给 \((1-r)/Z_0=t/Z_0\)，即 \(t=1-r\)。压力跃变 \(1+r-t=Z_s t/Z_0\)，所以 \(2r=(Z_s/Z_0)(1-r)\)，得到 \(\boxed{r=Z_s/(2Z_0+Z_s),\quad t=2Z_0/(2Z_0+Z_s)}\)。

(2) 两侧介质相同且特征阻抗为实数，反射率 \(\mathcal R=|r|^2\)、透射率 \(\mathcal T=|t|^2\)。因此 \(\mathcal A=1-\mathcal R-\mathcal T=4Z_0R/[(2Z_0+R)^2+X^2]\ge0\)。吸收功率还可由 \(\frac12 R|v|^2\) 求得，与散射能流差一致。

(3) 若 \(X=0\)，\(\mathcal A(R)=4Z_0R/(2Z_0+R)^2\)。求导为零在 \(R=2Z_0\)，极值为全局最大 \(\mathcal A_{\max}=1/2\)；此时 \(r=t=1/2\)，各占能流 \(1/4\)。对称单薄阻抗屏有两个无源辐射端口，单侧入射仍必须将部分能量分流到反射或透射端。

(4) \(Z_s\to0\) 时 \(r\to0,t\to1\)，屏几乎透明；\(|Z_s|\to\infty\) 时 \(r\to1,t\to0\)，近似刚性反射；\(R=0\) 时 \(\mathcal A=0\)，纯电抗只储存/释放近场能量而不产生净热耗散。
